\subsection{INVARIANTS OF A SEMI-SIMPLE LIE ALGEBRA RELATIVE TO A SEMI-SIMPLE ACTION}

In this no., k is assumed to be of characteristic zero.

PROPOSITION 13. Let g be a semi-simple Lie algebra, a a subalgebra of g reductive in g, and m the commutant of a in g. The subalgebra m satisfies conditions (1) and (2) of Lemma 2 of no. 3; it is reductive in g.

By Prop. 6 of Chap. I, §3, no. 5, applied to the a-module g, we have \(g = m \oplus [a, g]\) . Let B be the Killing form of g, and let \(x \in a, y \in m, z \in g\) . Then,

\[
\mathrm{B} ([ z, x ], y) = \mathrm{B} (z, [ x, y ]) = 0 \quad \text { since } [ x, y ] = 0,
\]

which shows that m is orthogonal to \([a,g]\) relative to B. Since B is non-degenerate, and since \(g = m \oplus [a,g]\) , this implies that the restriction of B to m is non-degenerate; condition (1) of Lemma 2 is thus satisfied.

Now let \(x \in m\) and let s and n be its semi-simple and nilpotent components. The semi-simple component of ad x is ad s, cf.~Chap. I, §6, no. 3. Since ad x is zero on a, so is ad s, by Prop. 4 (i). Thus \(s \in m\) , so \(n = x - s \in m\) , and condition (2) of Lemma 2 is satisfied.

Remark. The commutant of \(\mathfrak{m}\) in \(\mathfrak{g}\) does not necessarily reduce to \(\mathfrak{a}\) , cf.~Exerc. 4.

PROPOSITION 14. Let g be a semi-simple Lie algebra, A group and r a homomorphism from A to \(\text{Aut}(\mathfrak{g})\) . Let m be the subalgebra of g consisting of the elements invariant under \(r(A)\) . Assume that the linear representation r is semi-simple. Then m satisfies conditions (1) and (2) of Lemma 2 of no. 3; it is reductive in g.

The proof is analogous to that of the preceding proposition:

Let \(\mathfrak{g}^{+}\) be the vector subspace of \(\mathfrak{g}\) generated by the \(r(a)x - x\) , \(a \in \mathrm{A}\) , \(x \in \mathfrak{g}\) . The vector space \(\mathfrak{g}' = \mathfrak{m} + \mathfrak{g}^{+}\) is stable under \(r(\mathrm{A})\) . Let \(\mathfrak{n}\) be a complement of \(\mathfrak{g}'\) in \(\mathfrak{g}\) stable under \(r(\mathrm{A})\) . If \(x \in \mathfrak{n}\) , \(a \in \mathrm{A}\) , \(r(a)x - x \in \mathfrak{n} \cap \mathfrak{g}^{+} = 0\) , so \(x \in \mathfrak{m}\) and then \(x = 0\) since \(\mathfrak{m} \cap \mathfrak{n} = 0\) . Thus, \(\mathfrak{g} = \mathfrak{g}' = \mathfrak{m} + \mathfrak{g}^{+}\) . Let B be the Killing form of \(\mathfrak{g}\) and let \(y \in \mathfrak{m}\) , \(a \in \mathrm{A}\) , \(x \in \mathfrak{g}\) . Then

\[
\begin{array}{r l} \mathrm{B} (y, r (a) x - x) & = \mathrm{B} (y, r (a) x) - \mathrm{B} (y, x) \\ & = \mathrm{B} (r (a ^ {- 1}) y, x) - \mathrm{B} (y, x) \\ & = \mathrm{B} (y, x) - \mathrm{B} (y, x) = 0. \end{array}
\]

Thus m and \(g^{+}\) are orthogonal relative to B. It follows that the restriction of B to m is non-degenerate; hence condition (1) of Lemma 2. Condition (2) is immediate by transport of structure.

